\begin{figure}[h!]
	\centering
	\vspace{-0cm}
	\includegraphics[width = \textwidth]{../plots/0_images/Evo1}
	\caption{\textbf{Stellar Evolution of Sun-type star:} \textbf{a)} Gravitational collapse of a gas cloud forms a protostar which begin to fuse Hydrogen and arrived to the main sequence. \textbf{b)} After billions of years, star leaves the main sequence and become a Red Giant. At the tip of the Red Giant branch, star spontaneously fuses Helium with a short the episode of Helium flash and slowly moves to horizontal branch \textbf{c)} Intermediate star evolve through AGB phase, expanding upto the size of Jupiter orbit and loses its envelope to form a planetary nebula at the age of 12 billion years. \textbf{d)} Ejecting way the outer envelope, the remnant of dead star shines due to the degeneracy pressure of its electron and resides in lower left of HR diagram as a White Dwarf. Source: Sketch taken from \citep{kaufmann1978}}
	\label{stellar_evolution}
\end{figure}
